\documentclass[10pt,a4paper]{article}
\usepackage[margin=15mm]{geometry}
\usepackage{array,booktabs,tabularx,microtype}
\input{glyphtounicode}
\pdfgentounicode=1
\usepackage[hidelinks]{hyperref}
\hypersetup{pdftitle={Functional comparison of the two VSL tasks},
 pdfauthor={},pdfsubject={Task-first and peer-backed VSL formulations: traffic behavior and control}}
\pagestyle{empty}
\setlength{\parindent}{0pt}
\setlength{\parskip}{5pt}
\setlength{\tabcolsep}{5pt}
\renewcommand{\arraystretch}{1.08}
\newcolumntype{Y}{>{\raggedright\arraybackslash}X}
\begin{document}
\begin{center}
{\Large\bfseries Functional comparison of the two VSL tasks}
\end{center}
\vspace{-4pt}
\raggedright

Both tasks study a three-to-two lane drop with ending lane \textbf{E},
adjacent continuing lane \textbf{A} and far lane \textbf{B}. They use continuous
traffic mass, a finite initial cohort, one constant common speed advice,
and count all traffic's time until exit at the choke. The task-first
formulation uses a homogeneous approach and a dedicated merge/recovery
model; the peer-backed formulation uses lane-cell dynamics adapted from
Subraveti, Han and Carey. Their lane-information/control question is shared;
functional equivalence or a performance advantage is not assumed.

\begin{tabularx}{\linewidth}{@{}>{\raggedright\arraybackslash}p{0.16\linewidth}YY@{}}
\toprule
\textbf{Aspect} & \textbf{Original task-first formulation} & \textbf{Peer-backed formulation}\\
\midrule
\textbf{Traffic state and propagation}
& Homogeneous lane densities in the approach; a finite terminal uses cumulative arrivals,
a travel delay and point discharge (link-transmission-style).
& Lane densities in successive physical cells across the whole road
(cell-transmission-type). The current specification uses 100\,m cells and 6\,s steps.\\
\midrule
\textbf{Early merging}
& A fixed interaction length sets the E-to-A offer; spare A service limits
transfer. With sufficient A service, remaining E density decays exponentially
with distance traveled.
& Desired changes respond to perceived lane densities, downstream lookahead
and route/keep-right/cooperation incentives. Distance-to-end urgency remains;
shared service limits realized flow.\\
\midrule
\textbf{Permitted lane changes}
& Only E-to-A in preparation and at the choke. No ordinary lateral change
inside the terminal; B remains separate.
& Adjacent E\,$\leftrightarrow$\,A and A\,$\leftrightarrow$\,B changes are permitted
before the end, subject to incentives. At the end E must enter A or wait.\\
\midrule
\textbf{Advice and flow capacity}
& Advice changes the effective triangular flow--density diagram's critical
density and capacity. Undisturbed point-service ceilings stay fixed;
merge work and recovery constrain actual discharge.
& The nominal diagram stays fixed. Advice caps sending; Han's density-dependent
discharge/receiving laws govern congestion and capacity drop. Receiving
changes through traffic state.\\
\midrule
\textbf{Forced-merge allocation}
& Unqueued A gets priority if its arriving traffic fits. When both queue,
A/E share a work budget; E can require extra work, including a speed-mismatch cost.
& E's whole sending offer requests A service. Fixed equal priorities share
the receiver, capped by demand; unused shares are redistributed. Desire
does not raise priority. E has no extra work cost per unit.\\
\midrule
\textbf{Queues and storage}
& Its operating bounds keep the approach free, prevent terminal storage
blocking entry, and leave B unqueued. Waiting occurs at the A/E merging point.
& Finite storage and receiving constraints apply throughout the road.
Any lane can accumulate traffic; congestion can extend upstream across cells.\\
\midrule
\textbf{After-exit effects}
& Forced E entries retain an age- and insertion-speed-dependent influence
after completing their trips. Recent entries reduce service until their
downstream effect recovers.
& Exited traffic has no separate lingering restriction. Past movements affect
service through current lane-cell densities; no post-exit recovery account
is retained.\\
\midrule
\textbf{Travel-time objective}
& Complete travel time equals free travel plus A/E waiting. Merge delay is
generated by the service/recovery dynamics, then counted by total system time.
& Complete population-time is a finite trajectory sum plus an exact
remaining-time value. Total system time adds no special E-merging penalty.\\
\midrule
\textbf{Advice sensitivity}
& Continuous-time traffic dynamics and the timing of arrival and
queue-formation/clearance events must be differentiated together.
& Differentiate fixed-step stock updates and the exact remaining-time value
while active rules stay fixed; treat switching points separately.\\
\midrule
\textbf{Hidden lane information}
& E and total A+B can conceal an A/B imbalance. With homogeneous
initialization, adding initial A mass reconstructs the full initial state.
& E and A+B are observed at each section. One extra final-A density identifies
the two diagnostic states; broader fields may require more information.\\
\bottomrule
\end{tabularx}
\end{document}
